\chapter{调试这台机器}
\chaptermark{调试这台机器}
\label{sec:debugging}

\section{如何调试一台没有电路图的机器}

拿到一块能工作却没有电路图的电路板，工程师会用四种手法来调试它。他放上\emph{探针}，看导线上传的是什么。他\emph{注入测试信号}，看什么发生了变化。他\emph{断开}电路的一部分，看什么停止了工作。他还\emph{读取控制寄存器}，以了解电路板处于什么模式。整本书都是在尝试读出大脑的电路图；本章来看看，这四种手法中哪些工程师已经会用，以及前面各章对可以期待什么有何启示。

当今最引人注目的这类调试是脑机接口：它们读取神经元的脉冲并把它们变成外部机器的指令，或者反过来，从外部向大脑送入信号。本章大部分篇幅都讲它们，首先是 Neuralink 公司在做的事。

\section{探针：电极听到了什么}

大脑的探针是插入组织中的细电极。它能听到约一百微米半径内神经元的脉冲——通常是一个到几个。电极上的一个脉冲是几十到几百微伏、持续约一毫秒的尖峰，根据其波形可以把相邻的神经元区分开。最容易听到的是大的 \nt{GLUT} 神经元：它们在皮层中占多数（表~\ref{tab:types}），而且电流更强。

主要困难一眼可见。今天好的探针有几百或几千个电极，而大脑中有 $8.6\cdot10^{10}$ 个神经元。我们窃听的大约是这台机器的一亿分之一。为什么从如此微不足道的样本中竟能理解些什么？答案在第~\ref{sec:code}~章。相邻神经元的噪声是相关的，对群体求平均会碰到一个极限（命题~\ref{prop:pooling}）：一百个神经元携带的信息几乎和一千个一样多。此外，运动皮层要解决的问题是低维的：手臂的关节不多（第~\ref{sec:muscles}~章），控制它的几百个神经元的活动可以放进一个仅由一二十个公共变量构成的空间\cite{Gallego2017}。一个覆盖一百个神经元的探针几乎能听到所有这些变量。如果大脑在每个神经元里各存一条独立的消息，这样的窃听就毫无用处。

\section{数据流：直接在探针上压缩}

探针输出的数据流非常大。要看清脉冲的波形，每个电极需要以约十位的精度、每秒约 $20\,000$ 次进行数字化。一千个电极给出
\begin{empheq}[box=\keybox]{equation}
\begin{aligned}
R_{\text{原始}} \;&=\; N\,f_s\,b \;\approx\; 10^3\cdot2\cdot10^4\cdot10 \;\approx\; 2\cdot10^8\ \text{bit/s},\\
R_{\text{脉冲}} \;&\approx\; N\,r\,\log_2\frac{e}{r\,\Delta t} \;\approx\; 8\cdot10^4\ \text{bit/s}.
\end{aligned}
\label{eq:probedata}
\end{empheq}
第二个估计是 $r=10$ 次/秒、精度 $\Delta t=1\ms$ 时脉冲流的容量~\eqref{eq:spikeentropy}：其中真正包含的内容要小两千五百倍。对植入式设备来说这是决定性的差别：在头颅内能承受的、不会让组织发热的功耗下，原始数据流无法通过无线电穿过皮肤传出去。因此植入式设备需要在电极旁边的芯片上直接提取脉冲，只把事件——通道号和脉冲时刻——传出去。Neuralink 系统的第一份描述还没有做到这一点：芯片放大并数字化信号，整个原始数据流经电缆送出，脉冲由外部程序提取；芯片上的压缩和无线连接在那里被列为计划\cite{Musk2019}。

这是一个熟悉的手法。眼睛在把信号送上长电缆之前先将其压缩一百倍（第~\ref{sec:sensors}~章）；事件相机传输的不是帧，而是变化。为了挤进导线，探针不得不做大脑自己在做的事：用脉冲说话，而且只报告新消息。

\section{Neuralink 在做什么}

Neuralink 的设备就是为这些限制而打造的探针\cite{Musk2019}。它不用刚性的针阵列，而是用许多比头发还细的柔性细丝，每根细丝上都排列着电极：在系统的第一份描述中，最多 $3072$ 个电极分布在 $96$ 根细丝上，每根 $32$ 个；而在植入人体的设备中，据公司介绍，约一千个通道分布在 $64$ 根细丝上。柔性细丝对组织的损伤较小，也更能承受大脑在颅内不断轻微移动。细丝由机器人避开血管插入。如上所述，在第一份描述中原始数据流~\eqref{eq:probedata}~经电缆送出。据公司介绍，用于人体的设备结构不同：外壳完全被皮肤覆盖，脉冲在内部提取，数据经无线电送出，电力以无线方式输入。

第一台设备的任务是读取。细丝被放在计划手臂运动的那部分皮层中，解码器把脉冲变成屏幕上光标的运动。双手瘫痪的人通过想象运动来操作计算机。这个想法本身并不新：基于一百个电极刚性阵列的接口早已能让人控制光标\cite{Hochberg2006}、通过想象手写动作来写字\cite{Willett2021}，甚至能以每分钟约六十个词的速度把说话的尝试转换成屏幕上的文字\cite{Willett2023}。Neuralink 的新意在于工程：通道数、无线化、封闭的外壳和机器人插入，也就是尝试做一种可以在实验室外佩戴多年的探针。据我们所知，人体试验的结果主要由公司自己发布，而不是发表在同行评审的论文中；公司报告过，部分细丝在插入后发生了移位，工作的通道数因此减少。对调试来说这是常见的麻烦：相对电路板移动了的探针，听到的已经是别的导线。

公司的第二个方向是写入：一种为失明者向视觉皮层送入信号的视觉设备。下面在写入一节中再谈。

\section{读取：解码器作为接收方}

接口的解码器就是命题~\ref{prop:reader}~意义上的接收方：它自己决定读取消息的哪一部分。几乎所有接口读取的都是\emph{群体频率}：在几十毫秒的窗口内统计每个通道的脉冲数，再以一组权重相加，权重的选择使结果成为想要的运动。这就是第~\ref{sec:rate}~节中的群体频率编码，选择它并非偶然：每个窗口只有几个脉冲时，单个神经元的频率无法确定，而对一百个神经元求和就管用了。

能取出多少信息？用“意念之手”写字的速度约为每分钟九十个字符\cite{Willett2021}，即每秒一个半字符：即使每个字符五比特，也不到每秒八比特。据 Neuralink 介绍，控制光标可达每秒约十比特。而上界~\eqref{eq:spikeentropy}~对\emph{单个}神经元是每秒几十比特，对一百个神经元是几千比特。接口从被窃听的神经元中取出的，只是它们本可承载的一小部分。瓶颈不在导线，而在于群体携带多少个独立变量（命题~\ref{prop:pooling}），以及解码器对编码理解得有多好——正如第~\ref{sec:code}~章所见，这种编码尚未被完全破译。

还有第二个特点，我们在第~\ref{sec:motorlearning}~章已经熟悉。解码器和大脑相互学习。大脑看到光标去了哪里，得到误差并调整自己的指令——这正是经由误差导线的运动学习。而工程师不时重新拟合解码器的权重，因为细丝下的神经元在变化，网络中的连接在重新学习。于是有了两个相互适应的学习者；为了让这一对不失控，最好把它们的学习速度拉开：让一个学得快，另一个学得慢。

\begin{keyblock}
\begin{proposition}[为什么覆盖一千个神经元的探针能工作]
\label{prop:probe}
只窃听极小一部分神经元的接口之所以能控制运动，是因为群体活动是低维的，而且噪声是相关的：几百个神经元就携带了几乎全部公共变量。出于同样的原因，接口每秒只能取出几比特——其数量取决于任务中独立变量的个数，而不是脉冲流的容量。接口的工作性能已有测量；一旦理解了编码，解码器能改进多少仍是开放问题。
\end{proposition}
\end{keyblock}

\section{写入：比读取更难}

向机器送入信号比读取更难。通电的电极激发的不是一个神经元，而是它周围所有的神经元和导线，不分类型和符号。用它无法写入那种重要的是\emph{哪个}神经元、\emph{何时}发放的编码（第~\ref{sec:code}~章）：只能粗略地设定\emph{在哪里}以及\emph{有多强}。

对视觉来说，这在一定程度上够用了。人会把通过视觉皮层电极的电流看成视野中某个位置上的一个光点；当同时点亮多个光点、最多十六个时，失明的人能认出一些字母并分辨物体的边界\cite{Fernandez2021}。这是位置编码，它是可以写入的。但回想第~\ref{sec:sensors}~章：眼睛送进大脑的不是像素，而是压缩的描述——边缘、变化、变亮和变暗，分别走不同的导线。向皮层写入“光点”，就是向一台在输入端等待完全不同编码的机器写入像素。因此，人工视觉目前比按电极数目所能期待的要粗糙。也有一些不需要电极的测试信号：视错觉通过眼睛送入不寻常的输入，使机器偏离正常模式，而每一种错误都指向各自的机制（第~\ref{sec:illusions}~节）。

\section{探针看不到什么}

电极听到的是寻址总线——\nt{GLUT} 神经元的脉冲，以及较难听到的小 \nt{GABA} 神经元的脉冲。但在第~\ref{sec:serocomputer}--\ref{sec:acet}~章中我们看到，整台机器的工作模式由广播寄存器设定：\nt{SERO}、\nt{DOPA}、\nt{NORA}、\nt{ACET} 的浓度。电极听不到它们：源神经元少得可怜，而且它们的信号不是探针旁边的一个脉冲，而是一个体积中的浓度。能看到数据总线却看不到控制寄存器的调试者，总会感到意外：同一个输入会得到不同的回答，因为某处的 $g$、$a$ 或 $\tau_\gamma$ 变了。浓度可以用化学传感器直接在组织中测量\cite{Bunin1998}，但接口中还没有使用这类传感器。完全处于探针视野之外的还有机器的第二个、慢速的输出——血液中的激素（第~\ref{sec:chemoutput}~章）：它们是在血样中测量的，而不是用电极。

探针也看不到权重——而几乎所有记忆和学到的东西恰恰存储在权重中。只能根据回答如何变化来推测它们。探针也看不到人所感受到的疼痛：在第~\ref{sec:pain}~章中我们看到，警报信号的强度由机器自己设定——由脊髓中的闸门和上方的调节器设定，因此无法从传感器读出有多疼。

\section{小结：调试与开放问题}

\begin{table}[t]
\centering
\footnotesize
\setlength{\tabcolsep}{4pt}
\renewcommand{\arraystretch}{1.15}
\begin{tabular}{>{\raggedright}p{2.2cm}>{\raggedright}p{3.6cm}>{\raggedright}p{3.4cm}>{\raggedright\arraybackslash}p{4.0cm}}
\toprule
调试手法 & 接口做什么 & 本书的解释 & 已知什么 \\
\midrule
探针 & 听到几百到几千个神经元 & 低维性和相关噪声（第~\ref{sec:code}~章） & 已测量 \\
探针上的压缩 & 传输事件而不是原始信号 & 脉冲流的容量~\eqref{eq:spikeentropy} & 工程上已解决 \\
读取 & 群体频率 $\to$ 运动、文字、言语 & 解码器是接收方，群体频率编码 & 可用；每秒几比特 \\
共同学习 & 大脑和解码器相互适应 & 误差导线（第~\ref{sec:motorlearning}~章） & 已观察到；适应规则是研究课题 \\
写入 & 电流在视野中点亮光点 & 位置编码可以写入，时间编码不行（第~\ref{sec:sensors}~章） & 光点已测量；有用的视觉尚未实现 \\
读取寄存器 & 几乎没有 & 广播通道（第~\ref{sec:serocomputer}--\ref{sec:acet}~章） & 实验中已有化学传感器 \\
\bottomrule
\end{tabular}
\caption{调试这台机器：脑机接口能做什么，本书各章对它们有何解释。}
\label{tab:debugging}
\end{table}

表~\ref{tab:debugging}~汇总了本章。脑机接口已经能把探针放到寻址总线上，并从中每秒读出几比特；这件事之所以行得通，要归功于第~\ref{sec:code}~章的低维性和相关噪声。它们向机器写入的能力还很粗糙，因为机器的输入编码是压缩的，并且分别寻址到各条导线。而控制寄存器和权重——让机器可学习、可调节的东西——探针目前几乎看不到。

第~\ref{sec:synthesis}~章的每个开放问题同时也是调试者的任务。读取哪种编码，要等探针更密、解码器学会利用脉冲时间而不仅是频率时才能决定。多层网络如何学习，可以把探针同时放进几层，观察它们的回答在学习中如何变化来检验。广播通道传递的是什么，要等电探针加上化学探针后才能看清。本书尝试依据机器的行为和任何工程师都知道的定律来读出它的电路图。正在建造的调试工具将让我们能用探针检验这张电路图。

\section{习题}

\begin{exercise}[\lvlA]
\label{ex:17:1}
\emph{无线链路的预算。}穿过皮肤传输每比特需 $1$~nJ，发射机最多可用 $10$~mW。$N=1000$ 个电极时，原始数据流~\eqref{eq:probedata}~和脉冲流各需要多大功率？原始数据流要求每比特多少能量？
\end{exercise}

\begin{exercise}[\lvlA]
\label{ex:17:2}
\emph{一百个神经元有多少比特。}解码器在 $50\ms$ 窗口内对 $N$ 个神经元的脉冲求和，$r=20$~次/秒，噪声的两两相关 $c=0.1$，信号使群体频率改变 $30\%$（均方根）。估计 $N=10$、$100$、$1000$ 和 $N\to\infty$ 时每个窗口的速率 $\tfrac12\log_2(1+\text{SNR})$。$c=0$、$N=1000$ 时呢？
\end{exercise}

\begin{exercise}[\lvlB]
\label{ex:17:3}
\emph{键盘式接口的速率。}接口每秒一点五次从 $N=32$ 个字符中选一个：以概率 $P$ 选中想要的字符，错误等概率。$P=0.99$、$0.94$、$0.8$ 时它每秒传输多少比特？$P=0.94$ 时要达到 $10$~bit/s，每秒需要多少个字符？
\end{exercise}

\begin{exercise}[\lvlB]
\label{ex:17:4}
\emph{建模：群体解码器。}$N$ 个神经元，$\lambda_i=[b+m\,\mathbf v\cdot\mathbf p_i]_+$，$b=20$，$m=15$~次/秒，窗口 $50\ms$，$\mathbf v$ 每个分量的标准差为 $0.5$；所有神经元看到共同的噪声 $\mathbf v+\boldsymbol\xi$，$\sigma_\xi=0.2$。模拟无偏的群体向量解码器。$N$ 为多少时两种噪声相等？$N\to\infty$ 时每个窗口每个分量给出多少比特？
\end{exercise}

\begin{exercise}[\lvlB]
\label{ex:17:5}
\emph{两个学习者。}大脑输出指令 $m=b\,v^*$，解码器输出 $v=d\,m$，$\langle v^{*2}\rangle=1$。每次尝试后，两者都以速率 $\eta$ 沿 $\tfrac12(v-v^*)^2$ 做一步梯度下降。从 $b=1.2$、$d=1$ 出发，模拟 $\eta=0.8$、$1.1$、$1.5$、$2$。各自单独时在 $\eta<2$ 下稳定；这一对在 $\eta$ 为多少时稳定？
\end{exercise}

\begin{exercise}[\lvlB]
\label{ex:17:6}
\emph{设计：探针流水线。}$1024$ 个通道，每秒 $20\,000$ 个采样，神经元 $10$~次/秒，无线电 $1$~nJ/bit。对采样的白噪声取多高的阈值，才能使每个通道的虚假脉冲不超过 $0.1$~Hz？每 $20\ms$ 传输一次脉冲计数，每次 $2$ 比特：功率是多少？比传输原始 $10$ 位采样小多少倍？
\end{exercise}

\begin{exercise}[\lvlC]
\label{ex:17:7}
\emph{接口速率的极限。}对 $D=10$ 个带宽为 $W=5$~Hz 的变量，在 $\text{SNR}\approx0.9$（类似信号的噪声）和 $90$（$1000$ 个神经元的独立噪声）时估计 $R\le DW\log_2(1+\text{SNR})$。实测约为 $10$~bit/s。什么实验能区分这一差距的两种解释：类似信号的噪声和对编码的无知？
\end{exercise}
